\documentclass[master=cws,masteroption=gs]{kulemt}
\setup{title={Configuratieafhankelijkheden gebruiken om gedistribueerde applicaties effici\"ent te beheren in een hybride cloud},
  author={Harm De Weirdt},
  promotor={Prof.\,dr.\,ir.\ Wouter Joosen},
  assessor={Ir.\,W. Eetveel\and W. Eetrest}, %TODO
  assistant={Ir.\ Bart Vanbrabant}}
% De volgende \setup mag verwijderd worden als geen fiche gewenst is.
\setup{filingcard,
    translatedtitle={Configuratieafhankelijkheden gebruiken om gedistribueerde applicaties effici\"ent te beheren in een hybride cloud},
  udc=T134, %TODO: checken
  shortabstract={
      \begin{description}
          \item[Context] \hfill \\ Om IT infrastructuren efficient te beheren wordt er gebruik gemaakt van configuratiebeheergereedschappen. Deze gereedschappen zijn model gebaseerd, waarbij het model de gewenste toestand van de configuratie beschrijft. Om de configuratie door te voeren wordt de gewenste toestand vergeleken met de huidige toestand en worden de nodige acties afgeleid die nodig zijn om de infrastructuur in die gewenste toestand te brengen. Huidge systemen zijn reeds in staat om eenvoudige afhankelijkheden af te leiden. Bijvoorbeeld dat een service eerst geïnstalleerd moet worden voordat die service gestart kan worden. Wat ontbreekt is afhankelijkheden tussen services op verschillende systemen in rekening brengen.
        \item[Doel] \hfill \\Het doel van deze thesis is onderzoeken hoe afhankelijkheden in een configuratiemodel gebruikt kunnen worden om de initiele configuratie en mogelijke herconfiguraties van een hybrid cloud zo efficient mogelijk uit te voeren.
        \item[Onderzoeksvragen] \hfill \\
            \begin{enumerate}
                \item Hoe kunnen afhankelijkheden in een configuratiemodel gebruikt worden om veranderingen zo snel mogelijk uit te rollen?
                \item Kan de gevraagde tijd gesimuleerd worden in functie van het configuratie model?
            \end{enumerate}
        \item[Uitwerking] \hfill \\
            \begin{description}
                \item[Fase 1] Vertrouwd geraken met het configuratiebeheergereedschap
                \item[Fase 2] Onderzoeken van bestaande configuratiemodellen
                \item[Fase 3] Implementeren van een oplossing
                \item[Fase 4] Valideren van de oplossing door middel van de configuratiemodellen op een private en publieke cloud en een simulator
            \end{description}
    \end{description}
}}
% Verwijder de "%" op de volgende lijn als je de kaft wil afdrukken
%\setup{coverpageonly}
% Verwijder de "%" op de volgende lijn als je enkel de eerste pagina's wil
% afdrukken en de rest bv. via Word aanmaken.
%\setup{frontpagesonly}

% Kies de fonts voor de gewone tekst, bv. Latin Modern
\setup{font=lm}

% Hier kun je dan nog andere pakketten laden of eigen definities voorzien

% Tenslotte wordt hyperref gebruikt voor pdf bestanden.
% Dit mag verwijderd worden voor de af te drukken versie.
\usepackage[pdfusetitle,colorlinks,plainpages=false]{hyperref}

\usepackage{todonotes}

\begin{document}

\begin{preface}
  Dit is mijn dankwoord om iedereen te danken die mij bezig gehouden heeft.
  Hierbij dank ik mijn promotor, mijn begeleider en de voltallige jury.
  Ook mijn familie heeft mij erg gesteund natuurlijk.
\end{preface}

\tableofcontents*

\begin{abstract}
  In dit \texttt{abstract} environment wordt een al dan niet uitgebreide
  samenvatting van het werk gegeven. De bedoeling is wel dat dit tot
  1~bladzijde beperkt blijft.

\end{abstract}

% Een lijst van figuren en tabellen is optioneel
%\listoffigures
%\listoftables
% Bij een beperkt aantal figuren en tabellen gebruik je liever het volgende:
\listoffiguresandtables

\mainmatter

\chapter{Inleiding}
\label{inleiding}
Configuratiebeheergereedschappen laten toe op een effici\"ente manier IT infrastructuren op te zetten en onderhouden.
Daarvoor gebruiken ze een model dat de gewenste toestand van de volledige configuratie beschrijft.
De configuratie bestaat uit een reeks machines, elk met de bestanden, mappen, services,\ldots die ze moeten aanbieden.
Een eenvoudige voorbeeldconfiguratie is de LAMP-stack: 
\begin{itemize}
  \item Linux
  \item Apache webserver
  \item MySQL database (of MariaDB,MongoDB/\ldots)
  \item PHP (of Perl, Python,\ldots)
\end{itemize}

Bij het uitrollen van configuratie inspecteert de tool de huidige toestand van elke machine en vergelijkt ze met de gewenste
toestand.
Als er een verschil is maakt de tool de nodige aanpassingen. Indien niet onderneemt de tool geen actie.
Om de vergelijking uit te voeren moet eerst het model gecompileerd.
De huidige tools compileren voor elke machine hun deel van het model.
Dit zorgt ervoor dat bij de vergelijking de tool enkel informatie over die ene machine kan gebruiken.
Daardoor gebeurt het soms dat er teveel gecorrigeerd wordt en de gewenste toestand nog steeds niet bereikt is.
%TODO hoe vaak is soms? Wat triggert "soms"?
De toestand is wel al minder afwijkend van de gewenste situatie.
Een tool zal nooit aanpassingen maken die zorgen voor een configuratie die verder afwijkt van het model dan voorheen.
Na een paar iteraties zal dus altijd de gewenste configuratie bereikt worden.
%Ander woord voor tool? CMS? -> Definieren
Een grafische voorstelling van dit proces te zien op figuur \ref{fig:correctieAfwijking}.

\begin{figure}[p]
    \centering
    \includegraphics[width=0.8\textwidth]{images/correctieAfwijking.svg}
    \caption{Grafische voorstelling van de iteratief verkleinende afwijking van de gewenste situatie bij het uitrollen van een model.}
    \label{fig:correctieAfwijking}
\end{figure}

%TODO: de probleemstelling is nog niet echt gedefinieerd tot nu toe

%%% Local Variables: 
%%% mode: latex
%%% TeX-master: "masterproef"
%%% End: 

\chapter{Analyse van het probleem}
\label{chapter:1}
In dit deel wordt er dieper ingegaan op de probleemstelling.
Voor elke soort afhankelijkheid komt er een korte uitleg en daarna \'e\'en of meerdere heuristieken die toepasbaar zijn.

\section{Soorten afhankelijkheden}
\label{section:soorten_afhankelijkheden}
Zoals in de inleiding duidelijk werd bestaan er verschillende categori\"en van afhankelijkheden.
Elke categorie verschilt in moeilijkheid en aanpak. 

\subsection{Afhankelijkheden tussen bestanden en mappen}
\label{subs:bestanden_en_mappen}
#Hier gaat het over de afhankelijkheid tussen een bestand en de map waarin het staat.
#Een bestand kan niet bestaan zonder een map dus een heuristiek die elk bestand afhankelijk maakt van zijn parent folder is toepasselijk.
Tussen bestanden en mappen bestaat er zonder twijfel een sterke afhankelijkheid: een bestand kan helemaal niet bestaan zonder een map.
De eerste heuristiek die werd gebruikt in IMP is er een die voor elk bestand de "parent folder" opzoekt in het model en een afhankelijkheid toevoegd als deze gevonden wordt.
Als de map niet vermeld wordt in het model wordt deze ook niet toegevoegd omdat dit ongewenste gevolgen kan hebben. \todo{Welke?}
\missingfigure{Visuele voorstelling}

\subsection{Afhankelijkheden bij service stacks}
\label{subs:service_stacks}
Net zoals tussen bestanden en mappen is er een sterke afhankelijkheid tussen packages, services en hun eventuele configuratiebestanden:
een service kan niet gestart worden als zijn package niet ge\"installeerd is en zal niet correct werken zonder een juiste configuratie.
Een gelijkaardige heuristiek als die van bestanden en mappen kan er dus voor zorgen dat automatisch tussen de correcte afhankelijkheden worden ge\"introduceerd.
\\
"Service stack" is een eigen term die ik gegeven heb aan een verzameling van een package en een service (en mogelijks ook een configuratiebestand).
De verzameling wordt bepaald door te zoeken naar resources die binnen dezelfde scope gedefini\"eerd zijn.

\subsection{Afhankelijkheden door relaties}
\label{subs:relaties}
IMP laat toe in het model relaties tussen hoog-niveau objecten te specifi\"eren.
Sommige relaties, zoals "x [1:] -- [0:] y" impliceren dat y niet kan bestaan zonder een instantie van x.
Zo bekomen we een heuristiek die de afhankelijkheid tussen x en y toevoegt aan het model.

\subsection{Afhankelijkheden tussen machines}
In dit hoofdstuk wordt verder ingegaan op hoe afhankelijkheden tussen verschillende machines verwerkt worden.
Het belangrijkste punt is hoe hoog-niveau relaties en afhankelijkheden vertaalt worden naar dependencies op de lagere niveaus.

\section{Besluit van dit hoofdstuk}
 %Analyse van het probleem
\chapter{Een volgend hoofdstuk}
\label{hoofdstuk:2}
Een hoofdstuk behandelt een samenhangend geheel dat min of meer op zichzelf
staat. Het is dan ook logisch dat het begint met een inleiding, namelijk
het gedeelte van de tekst dat je nu aan het lezen bent.

\section{Eerste onderwerp in dit hoofdstuk}
De inleidende informatie van dit onderwerp.

\subsection{Een item}
Een tekst staat nooit alleen. Dit wil zeggen dat er zeker ook referenties
nodig zijn. Dit kan zowel naar on-line documenten\cite{wiki} als naar
boeken\cite{pratchett06:_good_omens}.

\section{Figuren}
Figuren worden gebruikt om illustraties toe te voegen. Dit is dan ook de
manier om beeldmateriaal toe te voegen zoals getoond wordt in
figuur~\ref{fig:logo}.

\begin{figure}
  \centering
  \includegraphics{logokul}
  \caption{Het KU~Leuven logo.}
  \label{fig:logo}
\end{figure}

\section{Tabellen}
Tabellen kunnen gebruikt worden om informatie op een overzichtelijke te
groeperen. Een tabel is echter geen rekenblad! Vergelijk maar eens
tabel~\ref{tab:verkeerd} en tabel~\ref{tab:juist}. Welke tabel vind jij het
duidelijkst?

\begin{table}
  \centering
  \begin{tabular}{||l|lr||} \hline
    gnats     & gram      & \$13.65 \\ \cline{2-3}
              & each      & .01 \\ \hline
    gnu       & stuffed   & 92.50 \\ \cline{1-1} \cline{3-3}
    emu       &           & 33.33 \\ \hline
    armadillo & frozen    & 8.99 \\ \hline
  \end{tabular}
  \caption{Een tabel zoals het niet moet.}
  \label{tab:verkeerd}
\end{table}

\begin{table}
  \centering
  \begin{tabular}{@{}llr@{}} \toprule
    \multicolumn{2}{c}{Item} \\ \cmidrule(r){1-2}
    Animal    & Description & Price (\$)\\ \midrule
    Gnat      & per gram    & 13.65 \\
              & each        & 0.01 \\
    Gnu       & stuffed     & 92.50 \\
    Emu       & stuffed     & 33.33 \\
    Armadillo & frozen      & 8.99 \\ \bottomrule
  \end{tabular}
  \caption{Een tabel zoals het beter is.}
  \label{tab:juist}
\end{table}

\section{Lorem ipsum}
Tenslotte gaan we hier nog wat tekst voorzien zodat er minstens een
bijkomende bladzijde aangemaakt wordt. Dat geeft de gelegenheid om eens te
zien hoe de koptekst en de voettekst zich gedragen.

\subsection{Lorem ipsum dolor sit amet, consectetur adipiscing elit}
Sed nec tortor id felis tristique sodales. Nulla nec massa eu dui fermentum
tincidunt. Integer ullamcorper ante eget eros posuere faucibus. Nam id
ligula ut augue pulvinar vulputate id at purus. Aenean condimentum tortor
eu mi placerat eget eleifend massa mollis. Nam est mi, sagittis quis
euismod eget, sagittis in nibh. Proin elit turpis, aliquam et imperdiet
sed, volutpat eu turpis.

Pellentesque vel enim tellus, vitae egestas turpis. Praesent malesuada elit
non nisi sollicitudin non blandit lacus tincidunt. Morbi blandit urna at
lectus ornare laoreet. Suspendisse turpis diam, lobortis dictum luctus
quis, commodo at lorem. Integer lacinia convallis ultricies. Sed quis augue
neque, eu malesuada arcu. Nullam vehicula, purus vitae sagittis pulvinar,
erat eros semper massa, eu egestas nibh erat quis magna. Cras pellentesque,
nisl eu dapibus volutpat, urna augue ornare quam, quis egestas lectus nulla
a lectus.

Vivamus dictum libero in massa cursus sed vulputate eros imperdiet. Donec
lacinia, libero ac lobortis egestas, nibh dui ornare arcu, luctus porttitor
velit massa sit amet quam. Maecenas scelerisque laoreet diam, vitae congue
quam adipiscing vitae. Aliquam cursus nisl a leo convallis eleifend
fermentum massa porta. Nunc libero quam, dapibus dapibus molestie sit amet,
faucibus vel nunc.

\subsection{Praesent auctor venenatis posuere}
Sed tellus augue, molestie in pulvinar lacinia, dapibus non ipsum. Fusce
vitae mi vitae enim ullamcorper hendrerit eu malesuada est. Proin iaculis
ante sed nibh tincidunt vel interdum libero posuere. Vivamus accumsan metus
quis felis congue suscipit dapibus enim mattis. Fusce mattis tortor eget
ipsum interdum sagittis auctor id metus.

Integer diam lacus, pharetra sit amet tempor et, tristique non lorem.
Aenean auctor, nisi eu interdum fermentum, lectus massa adipiscing elit,
sed facilisis orci odio a lectus. Proin mi nibh, tempus quis porta a,
viverra quis enim. In sollicitudin egestas libero, quis viverra velit
molestie eget. Nulla rhoncus, dolor a mollis vestibulum, lacus elit semper
nisi, nec sollicitudin sem urna eu magna. Nunc sed est urna, euismod congue
mi.

\subsection{Cras vulputate ultricies venenatis}
Vivamus eros urna, sodales accumsan semper vel, lobortis sit amet mauris.
Etiam condimentum eleifend lorem, ullamcorper ornare lectus aliquet vitae.
Praesent massa enim, interdum sit amet semper et, venenatis ut elit.
Quisque faucibus, quam ac lacinia imperdiet, nulla neque elementum purus,
tempus rutrum justo massa porta sapien. Vestibulum ante ipsum primis in
faucibus orci luctus et ultrices posuere cubilia Curae; Sed ultrices
interdum mi, et rhoncus sapien rutrum sed.

Duis elit orci, molestie quis sollicitudin sed, convallis non ante.
Maecenas tincidunt condimentum justo, et ultricies leo tristique vitae.
Vestibulum quis quam non lectus dapibus eleifend a vitae nibh. Nam nibh
justo, pharetra quis iaculis consequat, elementum quis justo. Etiam mollis
lacinia lacus, nec sollicitudin urna lobortis ac. Nulla facilisi.

Proin placerat risus eleifend erat ultricies placerat. Etiam rutrum magna
nec turpis euismod consectetur. Phasellus tortor odio, lacinia imperdiet
condimentum sed, faucibus commodo erat. Phasellus sed felis id ante
placerat ultrices. Aenean tempor justo in tortor volutpat eu auctor dolor
mollis. Aenean sit amet risus urna. Morbi viverra vehicula cursus.

\subsection{Donec nibh ante, consectetur et posuere id, tempus nec arcu}
Curabitur a tellus aliquet ipsum pellentesque scelerisque. Etiam congue,
risus et volutpat rutrum, est purus dapibus leo, non cursus metus felis
eget ligula. Vivamus facilisis tristique turpis, ut pretium lectus luctus
eleifend. Fusce magna sapien, ullamcorper vitae fringilla id, euismod quis
ante.

Phasellus volutpat, nunc et pharetra semper, sem justo adipiscing mauris,
id blandit magna quam et orci. Vestibulum a erat purus, ut molestie ante.
Vestibulum ante ipsum primis in faucibus orci luctus et ultrices posuere
cubilia Curae; Proin turpis diam, consequat ut ullamcorper ut, consequat eu
orci. Sed metus risus, fringilla nec interdum vel, interdum eu nunc.
Suspendisse vel sapien orci.

\subsection{Morbi et mauris tempus purus ornare vehicula}
Mauris sit amet diam quam, eget luctus purus. Sed faucibus, risus semper
eleifend iaculis, mi turpis bibendum nisl, quis cursus nibh nisl sit amet
ipsum. Vestibulum tempor urna vitae mi auctor malesuada eget non ligula.
Nullam convallis, diam vel ultrices auctor, eros eros egestas elit, sed
accumsan arcu tortor eget leo. Vestibulum orci purus, porttitor in pharetra
eget, tincidunt eget nisl. Nullam sit amet nulla dui, facilisis vestibulum
dui.

Donec faucibus facilisis mauris ac cursus. Duis rhoncus quam sed nisi
laoreet eu scelerisque massa tincidunt. Vivamus sit amet libero nec arcu
imperdiet tempor quis non libero. Sed consequat dignissim justo. Phasellus
ullamcorper, velit quis posuere vulputate, felis erat tincidunt mauris, at
vestibulum justo lectus et turpis. Maecenas lacinia convallis euismod.
Quisque egestas fermentum sapien eu dictum. Sed nec lacus in purus dictum
consequat quis vel nisl. Fusce non urna sem. Curabitur eu diam vitae elit
accumsan blandit. Nullam fermentum nunc et leo dictum laoreet. Donec semper
varius velit vel fringilla. Vivamus eu orci nunc.

\section{Besluit van dit hoofdstuk}
Als je in dit hoofdstuk tot belangrijke resultaten of besluiten gekomen
bent, dan is het ook logisch om het hoofdstuk af te ronden met een
overzicht ervan. Voor hoofdstukken zoals de inleiding en het
literatuuroverzicht is dit niet strikt nodig.

%%% Local Variables: 
%%% mode: latex
%%% TeX-master: "masterproef"
%%% End: 
 %Evaluatie
\chapter{Besluit}
\label{besluit}


% Indien er bijlagen zijn:
\appendixpage*          % indien gewenst
\appendix
\chapter{The First Appendix}
\label{app:A}
Appendices hold useful data which is not essential to understand the work
done in the master thesis. An example is a (program) source.
An appendix can also have sections as well as figures and references\cite{h2g2}.

\section{More Lorem}
\lipsum[50]

\subsection{Lorem 15--17}
\lipsum[15-17]

\subsection{Lorem 18--19}
\lipsum[18-19]

\section{Lorem 51}
\lipsum[51]

%%% Local Variables: 
%%% mode: latex
%%% TeX-master: "thesis"
%%% End: 


\backmatter
% Na de bijlagen plaatst men nog de bibliografie.
% Je kan de  standaard "abbrv" bibliografiestijl vervangen door een andere.
\bibliographystyle{abbrv}
\bibliography{referenties}

\end{document}

%%% Local Variables: 
%%% mode: latex
%%% TeX-master: t
%%% End: 
